\documentclass[12pt,a4paper]{article}

% Pakiety do obsługi języka polskiego i znaków
\usepackage[utf8]{inputenc}
\usepackage[T1]{fontenc}
\usepackage[polish]{babel}

% Pakiety matematyczne i graficzne
\usepackage{amsmath}
\usepackage{amssymb}
\usepackage{graphicx}
\usepackage{hyperref}
\usepackage{booktabs}
\usepackage{geometry}

% Ustawienia marginesów
\geometry{a4paper, margin=2.5cm}

\title{Dokumentacja Optymalizacji Modelu CNN \\ dla Zbioru Danych MNIST}
\author{Dokumentacja Techniczna PyTorch}
\date{\today}

\begin{document}

\maketitle
\tableofcontents
\newpage

\section{Wstęp}
Niniejsza dokumentacja techniczna opisuje proces modyfikacji i strojenia (ang. \textit{tuning}) Konwolucyjnej Sieci Neuronowej (CNN) służącej do klasyfikacji obrazów ze zbioru MNIST. Celem podjętych działań było zastosowanie technik zapobiegających przeuczeniu, wyjaśnienie anomalii funkcji straty, osiągnięcie skuteczności (ang. \textit{accuracy}) na poziomie powyżej 99\% na zbiorze testowym oraz wizualizacja pomyłek modelu za pomocą macierzy błędów.

\section{Zadanie 1: Regularyzacja L2 (Weight Decay)}
Aby zapobiec zjawisku przeuczenia w modelach głębokich (np. klasie \texttt{Deep}), wprowadzono regularyzację L2. W bibliotece PyTorch realizuje się to poprzez dodanie parametru \texttt{weight\_decay} do optymalizatora.

Regularyzacja L2 dodaje do standardowej funkcji straty $\mathcal{L}_0$ człon karzący za zbyt duże wagi, co można zapisać wzorem:
$$ \mathcal{L} = \mathcal{L}_0 + \lambda \sum_{i} w_i^2 $$
gdzie $\lambda$ to współczynnik siły regularyzacji (w naszym kodzie ustalony na np. $10^{-4}$), a $w_i$ to poszczególne wagi modelu. Ograniczenie wielkości wag zapobiega nadmiernemu dopasowywaniu się sieci do szumu obecnego w danych treningowych.

\section{Zadanie 2: Różnica między Validation Loss a Train Loss}
W początkowej fazie trenowania zauważono anomalię: strata na zbiorze walidacyjnym była niższa niż strata na zbiorze treningowym. Jest to bezpośredni wynik zastosowania warstwy \textbf{Dropout}.

Warstwa Dropout podczas treningu wyzerowuje losowo wybraną frakcję neuronów z prawdopodobieństwem $p$. Operacja ta sztucznie "upośledza" sieć na etapie uczenia, co podnosi wartość \textit{Train Loss}. 
Podczas ewaluacji (wywołanie \texttt{model.eval()}) na zbiorze walidacyjnym i testowym, mechanizm Dropout jest wyłączany. Model korzysta ze wszystkich wyuczonych wag jednocześnie, co skutkuje znacznie wyższą skutecznością i niższym błędem (\textit{Validation Loss}).

\section{Zadanie 3: Optymalizacja modelu CNN ($>99\%$ Accuracy)}
Aby przekroczyć próg 99\% poprawnych klasyfikacji na zbiorze testowym, wprowadzono następujące zmiany w architekturze \texttt{TunedCNN}:
\begin{itemize}
    \item \textbf{Zwiększenie liczby filtrów konwolucyjnych:} Pierwsza warstwa wyodrębnia 16 map cech (zamiast 2), a druga 32 (zamiast 4). Zwiększa to pojemność informacyjną modelu.
    \item \textbf{Zmiana klasyfikatora liniowego:} Warstwy w pełni połączone (Dense) zostały dostosowane do rozmiaru tensora po operacjach \texttt{MaxPool2d}. Liczba neuronów w warstwie ukrytej wynosi 128.
    \item \textbf{Zmiana optymalizatora:} Przejście ze stochastycznego spadku gradientu (SGD) na \textbf{Adam} (Adaptive Moment Estimation) ze współczynnikiem uczenia $\alpha = 0.001$, co znacząco przyspieszyło zbieżność.
\end{itemize}

\section{Zadanie 4: Macierz Błędów (Confusion Matrix)}
Do szczegółowej oceny jakości klasyfikacji poszczególnych cyfr wygenerowano macierz błędów $C$, w której każdy element $C_{i,j}$ jest zdefiniowany jako liczba obserwacji z prawdziwej klasy $i$, które zostały sklasyfikowane przez model do klasy $j$.

Dla problemu 10-klasowego (cyfry 0-9), diagonalna $C_{i,i}$ reprezentuje poprawnie rozpoznane próbki (\textit{True Positives}). Pozostałe elementy reprezentują pomyłki modelu. 
Przykładowo, element $C_{3,5}$ przechowuje informację, ile razy sieć zobaczyła na obrazku cyfrę "3", ale błędnie uznała ją za "5". Wizualizacja macierzy za pomocą biblioteki \texttt{scikit-learn} pozwala szybko zidentyfikować cyfry najczęściej mylone ze sobą, co zwykle dotyczy par $(4,9)$ oraz $(7,1)$ z powodu podobieństwa ich odręcznych kształtów.

\end{document}